\section{Introduction}\label{sec:introduction}
\PARstart{A}{n} evidence budget makes document retrieval a selection problem as well as a ranking problem. A system returning eight passages must decide whether a ninth candidate supplies useful additional information or merely repeats material already retrieved. Independent relevance estimates do not, by themselves, resolve this interaction. The distinction matters in retrieval-augmented generation, where selected passages occupy a limited context window and determine which evidence a generator can consult~\cite{lewis2020retrieval}. It also matters in ordinary document search: a compact list can contain individually plausible results while providing little variety.

Lower similarity between documents does not necessarily improve an answer. A passage that resembles another may contain the only correct support for a claim, while a dissimilar passage may be irrelevant. Relevance, redundancy, evidence coverage, and downstream answer quality must therefore remain separate quantities. The present work studies judged retrieval effectiveness and an embedding-based redundancy proxy. It neither substitutes that proxy for human evidence judgments nor attributes unmeasured generation outcomes to the retrieval stage.

Quantum search offers a useful setting in which to make these distinctions explicit. Grover search and amplitude amplification give an oracle-query improvement under their stated access assumptions~\cite{grover1996search,brassard2002amplitude}. A document pipeline, however, first constructs representations, retrieves candidates, computes relevance scores, and encodes an objective. A later amplitude-amplification step does not eliminate these operations. Simulating amplitudes on a classical computer adds work of its own. Its measured latency cannot be identified with the complexity of an ideal black-box quantum oracle.

Our starting point is the accepted Qubit Chat conference manuscript~\cite{conference2026}. That work reports controlled retrieval comparisons, an L6-guided Grover Adaptive Search (GAS) selector, a protected-prefix refinement, mathematical bounds, and a bounded native-oracle resource study. Those contributions and their experimental results are prior work. This journal article revisits the same completed evidence to answer questions that require finer analysis than a compact conference presentation permits: where relevant documents disappear from the candidate pipeline; how many queries improve or deteriorate despite a favorable mean; whether five stochastic seeds yield meaningfully different outcomes; and which measured stages account for the complete request cost.

The base policy selects eight documents from a hybrid shortlist of fifteen. A cross-encoder supplies the relevance component, while clipped pairwise embedding cosine measures redundancy. The subsequent refinement freezes the baseline's first five results, expands the candidate pool within the original dense top 64, and searches for a new three-document tail. Tail-score constraints restrict the search, and acceptance requires a strict decrease in the unquantized redundancy measure. Preserving the prefix fixes normalized discounted cumulative gain at rank five (nDCG@5), but it leaves recall and lower-ranked relevance free to change. This asymmetry is central to interpreting the evidence.

Three completed public test collections provide the empirical basis: SciFact, NFCorpus, and FiQA. SciFact and NFCorpus use the protected refinement; FiQA's reported row uses its earlier base policy. The article retains this distinction rather than combining metrics from different configurations into a single favorable row. The later failed FiQA refinement and NFCorpus confirmation exception remain part of the scientific record. Because earlier test outcomes informed method development, the follow-up tests support exploratory evaluation rather than an untouched confirmatory assessment.

The journal extension adds four analyses. We measure candidate-pool coverage and the best Recall@8 attainable within each saved pool, then examine query-level gains, losses, and ranking variation across seeds. We also reconstruct request-level latency to compare the costs of relevance scoring, base selection, and refinement. Finally, we connect these diagnostics to the recorded precision and resource grid and map the reported values to sealed artifacts. No new benchmark rankings, neural inference, or dataset collection are introduced.

The resulting claim is specific: constrained set selection can reduce measured redundancy while keeping a selected ranking prefix unchanged. Its recall consequences depend on the collection and on the relevance model's errors. Agreement with matched classical optimizers identifies a property of the formulation rather than a uniquely Grover-derived effectiveness gain. The resource evidence characterizes an implementation and its limits; it does not establish an operational quantum advantage.

The remainder of the article situates the formulation in retrieval and quantum optimization, specifies its algorithms and evaluation protocol, reports the recorded effectiveness results, and develops the journal-specific diagnostics. A separate resource section distinguishes compact simulation from native reversible-oracle synthesis. The discussion connects these findings to their validity limits and to experiments needed before stronger deployment claims could be made.
